% gen_event vs gen_server — Message Flow Comparison
% Visual comparison of 1-to-1 (gen_server) vs 1-to-N (gen_event) patterns.
% Compile with: xelatex gen_event_vs_gen_server.tex
\documentclass[border=10pt]{standalone}

% ============================================================
% Packages
% ============================================================
\usepackage{fontspec}
\usepackage{xeCJK}
\setCJKmainfont[BoldFont=STHeiti, ItalicFont=STKaiti]{STSong}
\setCJKsansfont{STHeiti}
\setCJKmonofont{STFangsong}

\usepackage{xcolor}
\usepackage{tikz}
\usepackage{pifont}

\usetikzlibrary{arrows.meta, positioning, shapes.geometric, fit, calc, backgrounds, decorations.pathreplacing}

% ============================================================
% Color Definitions
% ============================================================
\definecolor{erlred}{HTML}{A4070A}
\definecolor{erlblue}{HTML}{2B5EA7}
\definecolor{erlgreen}{HTML}{2E7D32}
\definecolor{erlyellow}{HTML}{F9A825}
\definecolor{erlpurple}{HTML}{6A1B9A}
\definecolor{erlmodule}{HTML}{D84315}
\definecolor{erlteal}{HTML}{00695C}
\definecolor{erlhandler}{HTML}{1565C0}
\definecolor{erlevent}{HTML}{E65100}
\definecolor{erlmanager}{HTML}{283593}
\definecolor{darkgray}{HTML}{333333}

% ============================================================
% Document
% ============================================================
\begin{document}

\begin{tikzpicture}[scale=0.72, every node/.style={scale=0.72},
    proc/.style={draw=#1!70, fill=#1!12, rounded corners=5pt,
                 minimum width=2.5cm, minimum height=1.2cm,
                 font=\small\bfseries, align=center},
    handler/.style={draw=erlhandler!70, fill=erlhandler!10, rounded corners=3pt,
                    minimum width=2cm, minimum height=0.8cm,
                    font=\tiny\ttfamily, align=center},
]

    % ==================================================================
    % TOP: gen_server model (1 Client -> 1 Server -> 1 Callback)
    % ==================================================================
    \node[font=\Large\bfseries\color{erlblue}] at (0, 8) {gen\_server：一对一模型};

    % Client
    \node[proc=erlpurple] (gs_client) at (-5, 6) {Client\\进程};

    % Server process
    \node[draw=erlblue!70, fill=erlblue!8, rounded corners=5pt,
          minimum width=5cm, minimum height=2.5cm] (gs_server) at (3, 6) {};
    \node[font=\small\bfseries\color{erlblue}, anchor=north] at (gs_server.north)
        {gen\_server 进程};

    % Single callback module inside
    \node[draw=erlgreen!70, fill=erlgreen!10, rounded corners=3pt,
          minimum width=4cm, minimum height=1.2cm,
          font=\scriptsize\ttfamily, align=center] (gs_mod) at (3, 5.6) {
        counter\_server\\
        State = \#\{count => 0\}
    };

    % Arrows
    \draw[-{Stealth[length=6pt]}, ultra thick, erlblue, line width=1.6pt]
        (gs_client.east) -- node[above, font=\scriptsize] {call / cast} (gs_server.west);
    \draw[-{Stealth[length=5pt]}, thick, erlred, line width=1.2pt]
        ([yshift=-3mm]gs_server.west) -- node[below, font=\scriptsize] {reply}
        ([yshift=-3mm]gs_client.east);

    % Label
    \node[font=\tiny\color{gray}, anchor=west] at (6.5, 6) {
        \begin{tabular}{l}
        1 进程\\
        1 回调模块\\
        1 份 State
        \end{tabular}
    };

    % ==================================================================
    % BOTTOM: gen_event model (N Sources -> 1 Manager -> N Handlers)
    % ==================================================================
    \node[font=\Large\bfseries\color{erlevent}] at (0, 3) {gen\_event：多对多模型};

    % Multiple event sources
    \node[proc=erlpurple] (src1) at (-6, 0.5) {Source A};
    \node[proc=erlpurple] (src2) at (-6, -1) {Source B};
    \node[proc=erlpurple] (src3) at (-6, -2.5) {Source C};

    % Event Manager process (large box)
    \node[draw=erlmanager!70, fill=erlmanager!5, rounded corners=5pt,
          minimum width=7cm, minimum height=5cm] (ge_mgr) at (4, -1) {};
    \node[font=\small\bfseries\color{erlmanager}, anchor=north] at (ge_mgr.north)
        {gen\_event Manager 进程};

    % Multiple handlers inside
    \node[handler] (h1) at (2.5, 0.2) {log\_handler\\State\_1};
    \node[handler] (h2) at (2.5, -1) {stats\_handler\\State\_2};
    \node[handler] (h3) at (2.5, -2.2) {alarm\_handler\\State\_3};

    % Callback details
    \node[font=\tiny\color{erlhandler!80!black}, anchor=west] at (4.2, 0.2) {handle\_event/2};
    \node[font=\tiny\color{erlhandler!80!black}, anchor=west] at (4.2, -1) {handle\_event/2};
    \node[font=\tiny\color{erlhandler!80!black}, anchor=west] at (4.2, -2.2) {handle\_event/2};

    % Arrows from sources to manager (converge)
    \coordinate (entry) at ([xshift=-5mm]ge_mgr.west);
    \draw[-{Stealth[length=5pt]}, thick, erlevent, line width=1.2pt]
        (src1.east) -- (entry |- src1.east) -- (entry);
    \draw[-{Stealth[length=5pt]}, thick, erlevent, line width=1.2pt]
        (src2.east) -- (entry);
    \draw[-{Stealth[length=5pt]}, thick, erlevent, line width=1.2pt]
        (src3.east) -- (entry |- src3.east) -- (entry);

    % notify label
    \node[font=\scriptsize\bfseries\color{erlevent}, anchor=south] at ([yshift=2pt]entry) {notify};

    % Fan-out arrows inside manager
    \draw[-{Stealth[length=4pt]}, thick, erlred!60, line width=0.8pt]
        ([xshift=2mm]ge_mgr.west) -- (h1.west);
    \draw[-{Stealth[length=4pt]}, thick, erlred!60, line width=0.8pt]
        ([xshift=2mm]ge_mgr.west) -- (h2.west);
    \draw[-{Stealth[length=4pt]}, thick, erlred!60, line width=0.8pt]
        ([xshift=2mm]ge_mgr.west) -- (h3.west);

    % Fan-out label
    \node[font=\tiny\bfseries\color{erlred!70}, rotate=90, anchor=south]
        at ([xshift=5mm]ge_mgr.west) {扇出};

    % Label
    \node[font=\tiny\color{gray}, anchor=west] at (7.5, -1) {
        \begin{tabular}{l}
        1 进程\\
        N 回调模块\\
        N 份独立 State
        \end{tabular}
    };

    % ==================================================================
    % DIVIDER
    % ==================================================================
    \draw[dashed, gray!50, line width=0.8pt] (-8, 3.8) -- (9, 3.8);

    % ==================================================================
    % BOTTOM COMPARISON TABLE
    % ==================================================================
    \node[draw=erlyellow!60, fill=erlyellow!5, rounded corners=4pt,
          minimum width=16cm, minimum height=2.8cm, anchor=north]
        (table) at (0.5, -4.5) {};
    \node[font=\small\bfseries\color{erlyellow!80!black}, anchor=north west]
        at ([xshift=4pt, yshift=-3pt]table.north west) {对比总结};

    \node[font=\tiny, text width=15cm, align=left, anchor=north west]
        at ([xshift=6pt, yshift=-16pt]table.north west) {
        \begin{tabular}{@{}l|l|l@{}}
        \textbf{维度} & \textbf{gen\_server} & \textbf{gen\_event} \\
        \hline
        进程模型 & 一个进程 = 一个服务 & 一个进程 = 一个事件管理器 \\
        回调模块 & 固定一个（启动时绑定） & 动态 N 个（运行时 add/delete） \\
        State & 一份，所有 handle\_* 共享 & N 份，每个 Handler 独立 \\
        消息分发 & 点对点（call/cast） & 广播（notify）+ 定向（call） \\
        容错 & 进程崩溃 = 服务不可用 & Handler 崩溃 = 仅移除该 Handler \\
        典型用途 & 有状态服务（DB连接、缓存） & 事件总线（日志、告警、统计） \\
        \end{tabular}
    };

\end{tikzpicture}

\end{document}
